\scanpage{306}
\mspage{47}
et pour \tikz[baseline=(q.base)]\node[draw,ellipse,inner sep=2pt](q){$i\in[0,N]$};, soit
\[
 \left\Vert\,\tau_i=\unsure{\eta|\upsilon}(t_i)\,.\right.
\]
Alors le type $\tau$ est donné par
\[
 \begin{aligned}
 \tau&=\tau_0\Delta_{\varepsilon_1}^{\deleted{\unsure{l_1|h_1}}}\tau_1
 \Delta_{\varepsilon_2}^{\deleted{\unsure{l_2|h_2}}}\ldots
 \Delta_{\varepsilon_N}^{\deleted{\unsure{l_N|h_N}}}\tau_N\\
 &\qquad\text{longueur}\qquad N+\sum_0^N\underset{\substack{\rule{0.3pt}{3pt}\rule{3.6em}{0.3pt}\rule{0.3pt}{3pt}\\\lambda_i\geq0}}{\operatorname{long}(\tau_i)}
 \end{aligned}\tag*{$(*)$}
\]
3\textsuperscript{o}) Description de \deleted{$\unsure{\mathrm{Fl}(\operatorname{Ob}|\operatorname{Ob}}$} $\mathrm{Fl}(\mathrm{Ty})$.

Si
\[
 f\colon\tau\longrightarrow\tau'\qquad\text{dans }\mathrm{Ty}
\]
on lui associe comme tantôt un entier $\uuline{N}\geq0$
($N=\operatorname{card}\mathrm{Cr}(\tau',f)-1$, \deleted{il est} $N=0$ ssi $\tau'=1$,
type de l'intervalle ponctuel), \uline{des} applications
\[
 \left\{\begin{aligned}
 [1,N]&\longrightarrow[+1,-1]\times\mathbb N^*,& i&\longmapsto(\varepsilon_i,\unsure{l_i|h_i})\\
 [0,N]&\longrightarrow\mathrm{Ty},& i&\longmapsto\tau_i\,.
 \end{aligned}\right.
\]
% slanted
\margin{\uline{NB} Si $N=0$ $[1,N]=\varnothing$ et il n'y a pas de $\varepsilon_i,l_i$ ; l'application
$[\underset{\substack{\parallel\\\{0\}}}{0,N}]\unsure{\mathbin{-}|\longrightarrow}\mathrm{Ty}$
se réduit à la donnée de $\tau_0\in\mathrm{Ty}$.}
La flèche $f$ dans $\mathrm{Ty}$ est déterminée par ces
données, le type $\tau$ est donné par la formule
$*$, ci-dessus, le type $\tau'$ par
\[
 \tau'=\Delta_{\varepsilon_1}^{l_1}\Delta_{\varepsilon_2}^{l_2}\ldots
 \Delta_{\varepsilon_N}^{l_N},\qquad\substack{\text{longueur}\\n'=\sum l_i}
 \tag*{$(*')$}
\]
et l'application $f$ croissante de $\inserted{I_\tau={}}[0,N+\sum\operatorname{long}\tau_i]$
dans $I_{\tau'}=[0,\sum l_i]$ est donnée ainsi
\[
 \begin{aligned}
 \deleted{\unsure{\Rightarrow}}\ \text{constante sur } &[0,0+\lambda_0],&\deleted{\unsure{\longmapsto}}\ 0\\
 \deleted{\unsure{\lambda_0+1|0+1}}\ \text{sur }&
 [\lambda_0+1,\lambda_0+\unsure{1|l}\deleted{\unsure{+1}}+\lambda_1]&\longmapsto l_1\\
 \text{\unsure{constant|constante} sur }&[\lambda_0+\lambda_1+\unsure{1|2},\lambda_0+\lambda_1+2+\lambda_2]&\unsure{\mathrel{\vdash\mkern-7mu\dashv}|\longmapsto}l_2\\
 \ldots&\\
 \text{\unsure{constante|constant} sur }&[\lambda_0+\lambda_{\unsure{1|i}}\deleted{\unsure{+}}\ldots+\lambda_i+(i+1),
 \lambda_0+\ldots+\lambda_i+\lambda_{i+1}+\unsure{(i+1)|(i-1)}]\unsure{\mathrel{\vdash\mkern-7mu\dashv}|\longmapsto}l_i
 \end{aligned}
\]
% archive sheet number: 305
